\chapter{Feature Engineering, Selection and Importance}\label{ch:ml:feature-engineering-selection-and-importance}

A model uses two versions of the same signal, computed over slightly different windows, whose correlation is 0.95. Its
\omterm{def:ml:feature-engineering-selection-and-importance:mdi}{permutation importance} ranks them fourth and fifth of fifteen. Refitted without the first, the model loses 0.09 points
of R-squared; without the second, it gains 0.02. Refitted without both, it loses 1.05 of its 2.67 points. Each copy looks
useless because the other stands in for it, and a researcher who pruned ``useless'' features one at a time would have
thrown away two-fifths of the model. This chapter builds features that behave, selects among them, and measures
importance with methods whose failures are known in advance.

\section{Transformations that respect the data}

\begin{definition}[Feature engineering]\label{def:ml:feature-engineering-selection-and-importance:engineering}
\emph{Feature engineering}\index{feature engineering} is the construction of model inputs from raw data: choosing what
to measure (a return, a volume, a book imbalance), over which window, and how to transform it (scale, rank, clip,
difference) so that its relation to the target is stable and its values are known at decision time.
\end{definition}

Market features are heavy-tailed and non-stationary. Book~7 (chapter~7) built the return, volatility and volume
features; Book~4 gave winsorisation (chapter~15) and fractional differencing (chapter~17). The transformation matters
more than the model for a feature with fat tails. A feature drawn from a Student $t$ with 1.5 degrees of freedom,
acting on the target through $0.2\tanh(z)$ (its effect saturates), explains 2.17\% of the variance at best. A linear
regression on the raw feature reaches 0.13\% out of sample: a few extreme values set the slope. Winsorised at the 1st
and 99th percentiles it reaches 1.27\%, ranked and mapped to normal scores 1.96\%. Ranking by date (the cross-sectional
rank transform of Book~7, chapter~6) is the default of this book for the same reason.

\section{Feature selection and stability selection}

\begin{definition}[Feature selection, stability selection]\label{def:ml:feature-engineering-selection-and-importance:selection}
\emph{Feature selection}\index{feature selection} chooses a subset of candidate features to use. \emph{Stability
selection}\index{stability selection} runs a selection procedure (the lasso's nonzero coefficients, a model's top
features) on many random half-samples and keeps the features selected in a share of them above a threshold
$\pi_{\mathrm{thr}}\in(\frac12, 1]$ (Meinshausen and Bühlmann, 2010).
\end{definition}

\begin{theorem}[A bound on false selections]\label{thm:ml:importance:mb}
If each half-sample selects on average $q$ of $p$ features, the selections of the noise features are exchangeable, and
the procedure is no worse than random guessing, the expected number of noise features selected by \omterm{def:ml:feature-engineering-selection-and-importance:selection}{stability selection} at
threshold $\pi_{\mathrm{thr}}$ is at most $q^2/((2\pi_{\mathrm{thr}} - 1)p)$.
\end{theorem}

\admitted{} The proof, by a Markov-type inequality on the selection frequencies, is in Meinshausen and Bühlmann (2010).

The chapter's task plants the roles: out of fifteen standard normal features, $x_1$ acts linearly, $x_2$ through its
square, $x_3$ and $x_4$ only through their product, $x_{1b}$ is $x_1$'s near-copy (correlation 0.95, no effect of its
own), and ten are noise; the signal is 4\% of the target's variance. For \omterm{def:ml:feature-engineering-selection-and-importance:selection}{stability selection} 85 more noise columns are
added (100 candidates) and 30 half-samples of 5\,000 rows are drawn. The lasso keeps $x_1$ alone at a threshold of 0.6
(it selects 6.7 features per half-sample; the bound allows 2.3 false selections, and none occurs); boosting's six
largest-gain features keep $x_1$ and $x_2$ (bound 1.8, none false). Neither finds $x_3$ and $x_4$: an interaction with
no main effects is invisible to a linear selector, and to a greedy tree whose first split on either feature gains
nothing. \omterm{def:ml:feature-engineering-selection-and-importance:selection}{Stability selection} controls what is selected, not what is missed.

\section{Importance: impurity, permutation, Shapley}

\begin{definition}[Mean decrease in impurity, permutation importance]\label{def:ml:feature-engineering-selection-and-importance:mdi}
The \emph{mean decrease in impurity}\index{mean decrease in impurity} (MDI, or gain importance) of a feature is the
total loss reduction of the splits on it across a tree ensemble, computed on the training data. The \emph{permutation
importance}\index{permutation importance} is the drop in a held-out score when the feature's column is randomly
permuted, breaking its link with the target and the other features.
\end{definition}

\begin{definition}[Shapley value, SHAP value]\label{def:ml:feature-engineering-selection-and-importance:shapley}
For a game $v$ on the set $P$ of players, the \emph{Shapley value}\index{Shapley value} of player $j$ is
\[
\phi_j = \sum_{S\subseteq P\setminus\{j\}}\frac{|S|!\,(|P| - |S| - 1)!}{|P|!}\bigl(v(S\cup\{j\}) - v(S)\bigr),
\]
the only allocation that is efficient ($\sum_j\phi_j = v(P) - v(\emptyset)$), symmetric, zero for a player that adds
nothing, and additive across games. A \emph{SHAP value}\index{SHAP value} is the Shapley value of feature $j$ for one
prediction, with $v(S)$ the model's expected output when the features in $S$ are fixed at their values (Lundberg and Lee,
2017); for tree ensembles it is computed exactly and fast (TreeSHAP).
\end{definition}

The four measures answer different questions. MDI says where the trees spent their splits on the training data;
\omterm{def:ml:feature-engineering-selection-and-importance:mdi}{permutation importance} says how much the fitted model needs a column on new data; drop-column importance (refit
without the feature) says how much the \emph{data} need it; the mean absolute \omterm{def:ml:feature-engineering-selection-and-importance:shapley}{SHAP value} says how much each feature moves
the model's individual predictions. On the planted task (\cref{fig:ml:importance:bars}), boosted trees fitted on 20\,000
rows score 2.67\% on 20\,000 others (the truth scores 4.01\%). Permutation and drop-column importance put $x_2$, $x_3$ and
$x_4$ first and give the noise almost nothing. MDI ranks the same four features first but hands every noise column a
share of the splits (the largest 5.4\% of the total, against 9.5\% for $x_1$): gain is earned on training data, noise
included. SHAP credits $x_{1b}$ more than $x_3$ or $x_4$: it explains what the model does, and the model uses the copy.

\begin{omfigure}
\begin{tikzpicture}
\begin{axis}[width=12cm,height=5.8cm,ybar,bar width=5.5pt,ylabel={share of the total (\%)},
  symbolic x coords={x1,x2,x3,x4,x1b,noise},xtick=data,xticklabels={$x_1$,$x_2$,$x_3$,$x_4$,$x_{1b}$,max noise},
  tick label style={font=\small},label style={font=\small},ymin=0,ymax=40,grid=major,grid style={gray!25},
  enlarge x limits=0.1,legend columns=4,legend style={font=\footnotesize,at={(0.5,-0.2)},anchor=north,draw=none,
  fill=none,/tikz/every even column/.append style={column sep=8pt}},table/col sep=comma]
\addplot[fill=gray!50,draw=gray] table[x=feature,y=gain]{figdata/ml/06-feature-engineering-selection-and-importance/importance.csv};
\addplot[fill=omDef!60,draw=omDef] table[x=feature,y=perm]{figdata/ml/06-feature-engineering-selection-and-importance/importance.csv};
\addplot[fill=omThm!55,draw=omThm] table[x=feature,y=drop]{figdata/ml/06-feature-engineering-selection-and-importance/importance.csv};
\addplot[fill=omProp!60,draw=omProp] table[x=feature,y=shap]{figdata/ml/06-feature-engineering-selection-and-importance/importance.csv};
\legend{gain (MDI),permutation,drop-column,mean $|$SHAP$|$}
\end{axis}
\end{tikzpicture}

\omcaption{Four importances of the same boosted trees on the planted task, each as a share of its positive total: the
four signal features, the copy $x_{1b}$, and the largest of the ten noise features. Permutation, drop-column and SHAP
are computed on the 20\,000 held-out rows, gain on the training rows. Data: \texttt{ml\_importance.importances}.}\label{fig:ml:importance:bars}
\end{omfigure}

\section{Substitution and clustered importance}

\begin{definition}[Substitution effect, clustered feature importance]\label{def:ml:feature-engineering-selection-and-importance:substitution}
The \emph{substitution effect}\index{substitution effect} is the understatement of a feature's importance when a
correlated feature can stand in for it: permuting or dropping one leaves the other. \emph{Clustered feature
importance}\index{clustered feature importance} measures importance for groups of correlated features, permuted or
dropped together, the groups found by hierarchical clustering on $1 - |\text{correlation}|$ (López de Prado, 2020).
\end{definition}

$x_1$ and its copy are the case. Permuted singly they cost 0.76 and 0.41 points of R-squared; permuted together, 2.72
points, more than the model's whole score, because the model has learned to use both and is now fed noise by both.
Dropped singly and refitted they cost 0.09 and $-0.02$ points; dropped together, 1.05. Clustering the fifteen features
at a distance of 0.5 finds exactly one group of two, $\{x_1, x_{1b}\}$, and every other feature alone: clustered
importance restores $x_1$'s signal to the pair and names the pair, which is the honest answer when the data cannot say
which copy matters. In MDI terms the pair also splits its credit (9.5\% and 8.1\%), so a pruning rule based on gain would
have hesitated over both.

\section{When importance lies}

\begin{method}[Reading importances]\label{met:ml:importance:reading}
\begin{enumerate}
\item Compute importances on held-out data (permutation, drop-column); read training-data MDI as a description of the
fit, never as evidence.
\item Cluster the features by correlation first and report the clusters' importances; prune whole clusters.
\item Compare with a permutation of the target (the canary of chapter~3): a feature is important only if it beats the
importance noise features receive.
\item Remember what each measure answers: the model's use (SHAP, permutation) or the data's need (drop-column); a
feature can be used and not needed.
\item Look for interactions separately: an effect through a product can be missed by every marginal screen.
\end{enumerate}
\end{method}

Importance also lies with time. A feature's importance in a model fitted on ten years is an average over regimes; a
feature that mattered in the first five years and not the last five gets half its old importance and looks
dependable. Chapter~12 measures importance on rolling windows, and chapter~21 asks whether an explanation is stable
across retrains.

\section{Tutorial: which feature mattered?}\label{tut:ml:feature-engineering-selection-and-importance:tut}

\begin{tutorial}
\textbf{Goal.} Compute the four importances and their clustered versions on the planted task, run \omterm{def:ml:feature-engineering-selection-and-importance:selection}{stability selection}
over 100 candidates, and transform a heavy-tailed feature three ways. \textbf{End state:} \cref{fig:ml:importance:bars},
the chapter's numbers.
\begin{tutsteps}
\item \textbf{Permutation and drop-column importance}, singly or by group.
\omcode{code/firm/featimp/firm_featimp.py}{48}{70}{Held-out importances, for single features or groups.}\label{lst:ml:importance:perm}
\item \textbf{Clusters and \omterm{def:ml:feature-engineering-selection-and-importance:selection}{stability selection}.}
\omcode{code/firm/featimp/firm_featimp.py}{77}{103}{Correlation clusters, stability selection and its bound.}\label{lst:ml:importance:stab}
\item \textbf{Run} \texttt{ml\_importance.importances()}, \texttt{twins()}, \texttt{stability()},
\texttt{heavy\_tail()} and \texttt{fig\_importance.py}.
\end{tutsteps}
\textbf{What to change next.} Give $x_3$ a small main effect and see whether \omterm{def:ml:feature-engineering-selection-and-importance:selection}{stability selection} now finds the
interaction; raise the copy's correlation to 0.999 and watch its \omterm{def:ml:feature-engineering-selection-and-importance:shapley}{SHAP value} overtake $x_1$'s.
\end{tutorial}

\section{Build: feature importance and selection}\label{bld:ml:feature-engineering-selection-and-importance:featimp}

\begin{build}
\bfield{Purpose} Every model report carries held-out, clustered importances and a record of how its features were
selected.
\bfield{Interface} \texttt{winsorise}, \texttt{rank\_gauss}, \texttt{mdi(model)},
\texttt{permutation\_importance(predict, X, y, score, reps, seed, groups)}, \texttt{drop\_column\_importance(fit\_predict,
X, y, Xt, yt, score, groups)}, \texttt{shap\_values(model, X)}, \texttt{cluster\_features(X, threshold)},
\texttt{stability\_selection(select, X, y, n\_sub, frac, seed)}, \texttt{false\_selection\_bound(q, p, threshold)}.
\bfield{Rules} Importances on held-out rows; groups permuted jointly with one permutation of the rows; the selection
procedure is a callable, so any model can be stability-selected.
\bfield{Acceptance tests} \texttt{code/firm/featimp/tests/}: a pure-noise feature has near-zero \omterm{def:ml:feature-engineering-selection-and-importance:mdi}{permutation importance}
and a duplicated signal is found only as a group; \omterm{def:ml:feature-engineering-selection-and-importance:shapley}{SHAP values} sum to the prediction minus the bias; clustering puts a
near-copy with its original; \omterm{def:ml:feature-engineering-selection-and-importance:selection}{stability selection} keeps a strong feature and drops noise; the bound's formula.
\bfield{Stretch} Conditional permutation (permute within bins of the correlated features); importance on rolling
windows; SHAP interaction values for pairs.
\end{build}

\begin{omsources}
\item N.~Meinshausen and P.~Bühlmann, ``Stability selection'', \emph{Journal of the Royal Statistical Society B} 72(4),
2010.
\item S.~M.~Lundberg and S.-I.~Lee, ``A unified approach to interpreting model predictions'', \emph{NeurIPS}, 2017;
S.~M.~Lundberg et al., ``From local explanations to global understanding with explainable AI for trees'', \emph{Nature
Machine Intelligence} 2, 2020.
\item C.~Strobl, A.-L.~Boulesteix, A.~Zeileis and T.~Hothorn, ``Bias in random forest variable importance measures'',
\emph{BMC Bioinformatics} 8, 2007.
\item M.~López de Prado, \emph{Machine Learning for Asset Managers}, Cambridge University Press, 2020.
\item L.~S.~Shapley, ``A value for $n$-person games'', in \emph{Contributions to the Theory of Games II}, 1953.
\end{omsources}

\section{Exercises}

\begin{exercise}[$\star$]\label{exo:ml:feature-engineering-selection-and-importance:1}
\omterm{def:ml:feature-engineering-selection-and-importance:selection}{Stability selection} over 200 candidates selects 10 per half-sample on average. What does the bound give at thresholds
0.6 and 0.9?
\end{exercise}

\begin{exercise}[$\star$]\label{exo:ml:feature-engineering-selection-and-importance:2}
Three players: $v(\emptyset) = 0$, singletons 1, 1, 0; pairs $\{1,2\} = 3$, $\{1,3\} = 1$, $\{2,3\} = 1$; all three 3.
Compute the \omterm{def:ml:feature-engineering-selection-and-importance:shapley}{Shapley values}.
\end{exercise}

\begin{exercise}[$\star$]\label{exo:ml:feature-engineering-selection-and-importance:3}
A held-out R-squared falls from 2.67\% to 2.58\% when a feature is dropped and the model refitted. What is its drop-column
importance, and what can you not conclude from it?
\end{exercise}

\begin{exercise}[$\star\star$]\label{exo:ml:feature-engineering-selection-and-importance:4}
Why does permuting $x_1$ and $x_{1b}$ together cost more (2.72 points) than the model's whole R-squared (2.67\%)?
\end{exercise}

\begin{exercise}[$\star\star$]\label{exo:ml:feature-engineering-selection-and-importance:5}
Why does MDI give the noise features up to 5.4\% of the total while \omterm{def:ml:feature-engineering-selection-and-importance:mdi}{permutation importance} gives them 0.2\%?
\end{exercise}

\begin{exercise}[$\star\star$]\label{exo:ml:feature-engineering-selection-and-importance:6}
\emph{Find the flaw.} ``We dropped every feature whose \omterm{def:ml:feature-engineering-selection-and-importance:mdi}{permutation importance} was below the median, one at a time,
refitting after each drop, and kept what was left.''
\end{exercise}

\begin{exercise}[$\star\star\star$]\label{exo:ml:feature-engineering-selection-and-importance:7}
\emph{Coding.} In \texttt{ml\_importance.heavy\_tail}, replace the $t$ with 1.5 degrees of freedom by one with 5. Report the
three R-squared and explain why the gap between raw and ranked closes.
\end{exercise}

\begin{exercise}[$\star\star\star$]\label{exo:ml:feature-engineering-selection-and-importance:8}
Show that if $x_{1b} = x_1$ exactly and a model's prediction depends on $x_1 + x_{1b}$ only, the \omterm{def:ml:feature-engineering-selection-and-importance:shapley}{SHAP values} of the
two are equal, and that the \omterm{def:ml:feature-engineering-selection-and-importance:mdi}{permutation importance} of the pair is not the sum of their single importances.
\end{exercise}

\section{Problem: Which Feature Mattered?}

\begin{problem}[{Weekend problem --- planted roles, four importances}]\label{pb:ml:feature-engineering-selection-and-importance:1}
The chapter's task: $x_1$ linear, $x_2$ squared, $x_3x_4$, a copy $x_{1b}$, ten noise features; boosted trees on 20\,000
rows, scored on 20\,000 more.

\medskip\noindent\textbf{Part I --- Transformations.}
\begin{enumerate}
\item What do the raw, winsorised and ranked versions of the heavy-tailed feature score, against the truth's 2.17\%?
\item Why does ranking beat winsorising here?
\item What does the book use by default, and why?
\item What does the model score, and the truth?
\end{enumerate}

\medskip\noindent\textbf{Part II --- Four importances.}
\begin{enumerate}[resume]
\item How does each measure rank $x_1$ to $x_4$?
\item Why does SHAP credit the copy more than $x_3$ or $x_4$?
\item What does MDI give the largest noise feature, and why?
\item Which measures answer the question ``does the data need this feature''?
\end{enumerate}

\medskip\noindent\textbf{Part III --- The copy.}
\begin{enumerate}[resume]
\item What do $x_1$ and $x_{1b}$ cost permuted singly and together?
\item What do they cost dropped singly and together?
\item Which group does the clustering find, at what distance?
\item What would one-at-a-time pruning have done?
\end{enumerate}

\medskip\noindent\textbf{Part IV --- Selection and the verdict.}
\begin{enumerate}[resume]
\item What do the lasso and boosting keep at a threshold of 0.6 over 100 candidates, and what does the bound allow?
\item Why do both miss $x_3$ and $x_4$?
\item State the \emph{named result}: the ranks of the signal features under each method, and the false-selection rate of
\omterm{def:ml:feature-engineering-selection-and-importance:selection}{stability selection} at 0.6.
\item Which importance would you show a portfolio manager, and which a risk committee?
\item How would you make \omterm{def:ml:feature-engineering-selection-and-importance:selection}{stability selection} see the interaction?
\item What would change with features that drift over time?
\item Write the importance section of a model report in four lines.
\item In one sentence: what does an importance measure?
\end{enumerate}
\end{problem}

\section{Interview questions}

\begin{interviewq}[$\star$ \iqroles{researcher, mle}]\label{iq:ml:feature-engineering-selection-and-importance:1}
What is the difference between impurity-based and \omterm{def:ml:feature-engineering-selection-and-importance:mdi}{permutation importance}, and which do you trust?
\end{interviewq}

\begin{interviewq}[$\star\star$ \iqroles{researcher}]\label{iq:ml:feature-engineering-selection-and-importance:2}
Two of your features are 95\% correlated and both show low importance. What do you do?
\end{interviewq}

\begin{interviewq}[$\star\star$ \iqroles{mle}]\label{iq:ml:feature-engineering-selection-and-importance:3}
Explain \omterm{def:ml:feature-engineering-selection-and-importance:shapley}{SHAP values} and one situation where they mislead.
\end{interviewq}

\begin{interviewq}[$\star\star$ \iqroles{researcher}]\label{iq:ml:feature-engineering-selection-and-importance:4}
How does \omterm{def:ml:feature-engineering-selection-and-importance:selection}{stability selection} control false discoveries, and what does it not control?
\end{interviewq}

\begin{interviewq}[$\star\star$ \iqroles{researcher, trader}]\label{iq:ml:feature-engineering-selection-and-importance:5}
A volume feature has occasional values a hundred times its median. How do you feed it to a model?
\end{interviewq}

\begin{interviewq}[$\star\star\star$ \iqroles{researcher}]\label{iq:ml:feature-engineering-selection-and-importance:6}
Construct a target and two features such that each feature alone has zero correlation with the target but together they
predict it perfectly. What does this imply for feature screening?
\end{interviewq}
